
过去十年的诸多发现为\textbf{\trHamiltonianSystem}（Hamiltonian system）这一古老领域开启了新的视野，并催生了一个新的领域：\textbf{\trSymplecticTopology}（symplectic topology）。
令人惊奇的\textbf{\trRigidity}（rigidity）现象表明，\textbf{\trSymplectic\trMapping}（symplectic mapping）的本性迥异于\trVolumePreserving{}映射（volume-preserving mapping），由此引出了许多新问题，其中不少至今悬而未决。
另一方面，由于对\textbf{\trClassicalMechanics}（classical mechanics）中一个古老\textbf{\trVariationalPrinciple}（variational principle）的分析，\trHamiltonianSystem{}中全局的\trPeriodic{}现象得以确立。
事实证明，这些看似不同的现象神秘地相互关联着，而联结之一便是一类\textbf{\trSymplecticInvariant}（symplectic invariant），即所谓的\textbf{\trSymplecticCapacity}（symplectic capacity）。
这些\trInvariant{}量正是本书的主题；本书源自作者们1991年在罗格斯大学（Rutgers University）、波鸿鲁尔大学（RUB Bochum）与苏黎世联邦理工学院（ETH Z\"urich），以及1992年在伯尔尼的\trBorel{}讨论班上所授的课程。
由于这些课程不要求任何预备知识，书中只选取了少数相当初等的专题并详细加以证明。
此外，我们的选取出于一条单一原理的指引：力学的\textbf{\trAction}（action）原理。
\textbf{\trPhaseSpace}（phase space）中\textbf{\trLoop}（loop）的\textbf{\trActionFunctional}（action functional）
\begin{equation*}
F(\gamma)=\int_\gamma p\,dq-\int_0^1 H\bigl(t,\gamma(t)\bigr)\,dt,
\end{equation*}
不同于\textbf{\trConfigurationSpace}（configuration space）中由一个\trLagrange{}量所定义的古老\trHamiltonian{}原理。
$F$ 的\textbf{\trCriticalPoint}（critical point）是那些求解与\textbf{\trHamiltonianFunction}（Hamiltonian function） $H$ 相关联的\textbf{\trHamiltonianEquation}（Hamiltonian equation）的\trLoop{} $\gamma$，因而也就是\textbf{\trPeriodicOrbit}（periodic orbit）。
这一\trVariationalPrinciple{}有时被称为\textbf{\trLeastActionPrinciple}（least action principle）。
然而，$F$ 没有极小值。
事实上，\trAction{}原理是非常\trDegenerate{}的：它的所有\trCriticalPoint{}都是\trMorseIndex{}为无穷的\textbf{\trSaddlePoint}（saddle point），而且乍看之下，该原理对存在性证明显得全然无用。
但出人意料，它非常有效。
本书将用若干\trVariational{}技巧演示这一点：始于P.~\trRabinowitz{}的\textbf{\trMinimax}（minimax）论证，而终于A.~\trFloer{}的同调。
本书包含如下论题：

导论性的一章以相当不系统的方式介绍了若干背景材料。
我们给出\textbf{\trSymplecticManifold}（symplectic manifold）与\trSymplecticMapping{}的定义，并简要回顾\textbf{\trHamiltonianFormalism}（Hamiltonian formalism）。
为方便起见，书中使用了\trCartanCalculus。
书中证明了：二维\trSymplecticManifold{}可按\textbf{\trEulerCharacteristic}（Euler characteristic）与总\textbf{\trVolume}（volume）分类。
一些后文将详细处理的问题在这里通过专门的例子提出并加以讨论。
为了在 $\R^{2n}$ 中\trHamiltonianSystem{}的一个\textbf{\trConvex\trEnergySurface}（convex energy surface）上建立\trPeriodicOrbit，我们演示了\textbf{\trCalculusOfVariations}（calculus of variations）的所谓\textbf{\trDirectMethod}（direct method）。
引入\trBirkhoff\trInvariant{}量，是为了在不给出证明的情况下刻画\trHamiltonianSystem{}在一个\textbf{\trEquilibrium}（equilibrium）附近或一个\textbf{\trPeriodicSolution}（periodic solution）附近错综复杂的\textbf{\trOrbit}（orbit）结构。
这些局部结果与后继各章所处理的全局问题形成了鲜明对比。

第2章以公理化的方式系统地引入了称为\trSymplecticCapacity{}的\trSymplecticInvariant。
考察固定维数 $2n$ 的\trSymplecticManifold{}全体所成之族，一个\textbf{\trCapacity}（capacity） $c$ 是给每个\trSymplecticManifold{} $(M,\omega)$ 指定一个正数 $c(M,\omega)$ 或 $\infty$ 的\trMapping，它满足以下公理：关于\textbf{\trSymplecticEmbedding}（symplectic embedding）的\textbf{\trMonotonicity}（monotonicity）公理、关于\textbf{\trSymplecticStructure}（symplectic structure）的\textbf{\trConformality}（conformality）公理，以及一条排除高维\trVolume{}的\textbf{\trNormalization}（normalization）公理。
对 $\R^{2n}$ 的子集而言，\trCapacity{}把\textbf{正\trQuadraticForm}（positive quadratic form）的人所熟知的\trLinear\trSymplecticInvariant{}（linear symplectic invariant）延拓到了\trNonlinear\trSymplecticMapping{}上。
若 $M$ 与 $N$ \trSymplectic\trDiffeomorphism{}，则 $c(M,\omega)=c(N,\tau)$。
鉴于其\trMonotonicity，一个\trCapacity{}尤其构成对某些\trSymplecticEmbedding{}的\textbf{\trObstruction}（obstruction），我们将用它来解释由Ya.~\trEliashberg{}与M.~\trGromov{}发现的\trSymplecticEmbedding{}的\trRigidity{}现象。
特别是，\trGromov{}的\trSqueezingTheorem{}将用\trCapacity{}推出，\trEliashberg{}的\textbf{\trSymplectic\trDiffeomorphism}（symplectic diffeomorphism）的 $C^0$ 稳定性亦然。
我们引入了\textbf{\trSymplecticHomeomorphism}（symplectic homeomorphism）的概念，它引出了许多问题。
\trCapacity{}函数有很多种。
例如，能够\trSymplectic{}地\trEmbedding{}一个\trSymplecticManifold{} $(M,\omega)$ 的 $\R^{2n}$ 中最大\textbf{\trBall}（ball）的尺寸，给出一种特殊的\trCapacity{}，称为\trGromovWidth。
它是最小的\trCapacity{}函数，最初由M.~\trGromov{}引入。
此外还有许多其他的「嵌入」\trCapacity{}。

第3章完全致力于对一个特定\trSymplecticCapacity{} $c_0$ 的非常详尽的构造。
它由\trHamiltonianSystem{}以动力学方式定义。
它度量\trHamiltonianFunction{} $H\colon M\to\R$ 的最小 $C^{0}$\textbf{\trOscillation}（oscillation）——即使得可断言相应 $X_H$ 在 $M$ 上存在快速周期解的振荡。
在\trConnected{}的二维\trSymplecticManifold{}这一特殊情形，\trCapacity{} $c_0$ 与总面积一致。
存在性证明基于上述\trAction{}原理，该原理将在其恰当的\textbf{\trFunctional{}分析}（functional analysis）框架中从头引入。
这一原理的有趣之处在于它既不\trBoundedBelow{}也不\trBoundedAbove，致使\trStandard{}的\trVariational{}技巧无法直接应用。
溯源于P.~\trRabinowitz{}的技巧使我们能够有效地确立该\trFunctional{}的特定\trSaddlePoint{}，它们代表着系统的特殊\trPeriodicSolution。
在\textbf{\trBounded{}\trSmooth\trConvex\trDomain}（bounded smooth convex domain） $U\subset\R^{2n}$ 的特殊情形，\trCapacity{}由其\textbf{\trBoundary}（boundary） $\partial U$ 的一条特定\textbf{\trClosedCharacteristic}（closed characteristic）体现：它具有等于 $c_0(U)$ 的极小\textbf{\trReducedAction}（reduced action）。
但一般说来，计算\trInvariant{}量 $c_0$ 相当困难。
近来某些基于一阶\trElliptic{}系统与\trFredholm{}理论的更高深技术的计算，书中不加证明地作了介绍。
随着\trCapacity{} $c_0$ 的构造，第2章所述\trRigidity{}现象的证明即告完备。
得益于其特殊性质，这个\trInvariant{}量还将证明对\trHamiltonianSystem{}的动力学颇为有用。

第4章把动力学\trCapacity{} $c_0$ 应用于\trHamiltonianSystem{}定性理论中一个源自\textbf{\trCelestialMechanics}（celestial mechanics）的古老问题：\textbf{\trCompact\trEnergySurface}（compact energy surface）是否承载\trPeriodicOrbit？
我们将证明，许多先前经由技术上繁复的证明得到的著名全局存在性结果，都能从这个\trInvariant{}量立即得出。
现象其实很简单：若\trSymplecticManifold{}中的一个\textbf{\trCompact\trHypersurface}（compact hypersurface）拥有\trCapacity{} $c_0$ 有限的\textbf{\trNeighborhood}（neighborhood），那么它的附近总有不可数多条\trClosedCharacteristic。
如果再附加在辛意义下\trInvariant{}的限制条件，例如「\textbf{\trContactType}（contact type）」，那么该\trHypersurface{}自身便承载一条\trClosedCharacteristic。
我们将特别证明由C.~\trViterbo{}给出的 $\R^{2n}$ 中\trWeinstein{}猜想这一开创性的解决。
一个非\trStandard{}的\textbf{\trSymplectic\trTorus}（symplectic torus）表明，与上文提到的\trGromovWidth{}相反，并非每个\trCompact\trSymplecticManifold{}都具有有限的\trCapacity{} $c_0$。
我们的特殊例子与M.~\trHerman{}关于\textbf{\trClosingLemma}（closing lemma）的著名\trCounterexample{}有关，该\trCounterexample{}回答了\textbf{\trDynamicalSystem}（dynamical system）中一个长期悬而未决的问题。
M.~\trHerman{}的「非闭合引理」（non-closing-lemma）在本章末尾给出证明。

第5章研究 $\R^{2n}$ 的\trSymplectic\trDiffeomorphism{}中由\trCompactSupport{}的时变\textbf{\trHamiltonianVectorField}（Hamiltonian vector field）生成的子群 $\mathcal{D}$。
这种\trSymplectic\trDiffeomorphism{} $\varphi$ 到\textbf{\trIdentityMap}（identity map）的\textbf{\trDistance}（distance）——或者说它的\textbf{\trEnergy}（energy） $E(\varphi)$——将借助其生成\trHamiltonianFunction{}的\trOscillation{}来度量。
由此得到 $\mathcal{D}$ 上一种出人意料的\textbf{\trBiInvariantMetric}（bi-invariant metric），称为\trHofer{}度量，定义为 $d(\varphi,\psi)=E(\varphi^{-1}\circ\psi)$。
该定义不涉及\trHamiltonianFunction{}的\textbf{\trDerivative}（derivative），具有 $C^0$ 本性。
验证\textbf{\trMetric}（metric）性质——即 $d(\varphi,\psi)=0$ 当且仅当 $\varphi=\psi$——是困难所在。
它基于对\trActionFunctional{}的更精细的\trMinimax{}论证，这些论证同时对一大类\trHamiltonianFunction{}成立。
我们将研究这一特殊度量与第3章引入的动力学\trSymplecticInvariant{} $c_0$ 的关系，以及它与另一个定义在 $\R^{2n}$ 子集上、称为\textbf{\trDisplacementEnergy}（displacement energy）的\trSymplecticInvariant{}的关系。
子集 $U$ 的\trDisplacementEnergy{}度量的是把给定集合 $U$ 从其自身挪开所需的最小\trEnergy{} $E(\varphi)$，这里挪开即指 $U\cap\varphi(U)=\varnothing$。
\trBiInvariantMetric{}还将与\trStandard{}的上确界度量（sup-metric）作比较。
与该\trMetric{}相联系的\textbf{\trGeodesic}（geodesic）弧将得到定义与详细的描述。
\trGeodesic{}弧的一个特殊例子是由一个自治\trHamiltonianFunction{}生成的\textbf{\trFlow}（flow）。
在我们的进路中，\textbf{\trHamiltonianMap}（Hamiltonian map） $\varphi\in\mathcal{D}$ 的\textbf{\trActionSpectrum}（action spectrum）扮演着重要角色，事实证明，它是实数的一个\trNowhereDense{}子集。
我们的\trMinimax{}原理从 $\mathcal{D}$ 上的\textbf{\trActionSpectrum{}\trBundle}（action spectrum bundle）中挑出一个非平凡的连续\textbf{\trSection}（section），称为 $\gamma$-\trInvariant{}量。
这个\trInvariant{}量是本章的主要技术工具。
它使我们能够刻画\trGeodesic，也用于证明\trCompactSupport{}的\trHamiltonianMap{}具有无穷多个非平凡\trPeriodic{}点的存在性。

第6章的主题是\trCompact\trSymplecticManifold{} $(M,\omega)$ 上\trHamiltonianMap{}的\textbf{\trFixedPoint}（fixed point）理论。
它不同于\trTopology{}学中的\trFixedPoint{}理论。
\trHamiltonianMap{}是一类特殊的\trSymplecticMapping：它与\trIdentityMap\trHomotopic{}，而且该\textbf{\trHomotopy}（homotopy）由一个时变\trHamiltonianVectorField{}的\trFlow{}生成。
受H.~\trPoincare{}的最后几何定理的推动，V.~I.~\trArnold{}在二十世纪六十年代猜想：这类\trHamiltonianMap{}具有\trFixedPoint{}的数目，至少与 $M$ 上实值函数具有\trCriticalPoint{}的数目一样多。
用\trDynamicalSystem{}的语言重新表述，这个猜想要求一种\trLjusternik-\trSchnirelman{}理论，或者对应地，要求一种关于\textbf{\trForcedOscillation}（forced oscillation）的\trMorseTheory{}，这些\trForcedOscillation{}是 $M$ 上时间周期\trHamiltonianSystem{}的解。
我们将先就\trStandard\trTorus{}（standard torus） $T^{2n}$ 这一特殊情形证明该猜想。
证明仍基于\trAction{}原理，但这一次目标在于找出它的所有\trCriticalPoint。
我们的策略受到C.~\trConley{}的\trDynamicalSystem{}\trTopology{}进路的启发：\textbf{\trRegularization\trGradient{}方程}（regularized gradient equation）正是由定义在\textbf{\trManifold}（manifold） $M$ 的\trContractible\trLoop{}集合之上的那个\trActionFunctional{}所产生的，我们将研究其全体\trBounded{}解所成集合的\textbf{\trTopology}（topology）。
这样一来，无穷维\textbf{\trLoopSpace}（loop space）中\trGradient{}流的研究便被化归为一个\trCompact\trMetric{}空间的\trGradientLike{}连续\trFlow{}的研究，后者的\textbf{\trRestPoint}（rest point）正是所需的\trCriticalPoint。
它们的个数进而用6.3中介绍的\trLjusternik-\trSchnirelman{}理论来估计。
经过重新解释，这将引导我们证明\trArnold{}猜想对满足 $[\omega]|\pi_2(M)=0$ 的更大一类\trSymplecticManifold{}成立。
在这种一般情形下没有自然的\textbf{\trRegularization}（regularization）可用，我们不得不在6.4中研究非\trRegularization\trGradient{}系统的有界解集合，这些解此时是一个特殊的一阶\trElliptic{}\trPartialDifferentialEquation{}组（\trCauchyRiemann{}型）的光滑解。
这些解与M.~\trGromov{}在 $M$ 中的\trPseudoholomorphic{}曲线有关。
解集合的\textbf{\trCompactness}（compactness）将基于一种分析技术，它有时被称为\textbf{\trBubblingOffAnalysis}（bubbling-off analysis）。
循此路线，我们将抵达这些发展的顶峰：A.~\trFloer{}对\trMorseTheory{}的全新进路与\trFloer{}同调。
对\trFloer{}美妙的思想，我们仅在6.5中作一概述。
\trFloer{}的进路与动力学\trCapacity{} $c_0$ 的构造相结合，产生了一种\textbf{\trSymplectic\trHomology}（symplectic homology）理论，它尚未定型，我们将在本章最后一节中不加证明地勾画。
这些理论在技术上要求的预备相当高深，超出了本书的范围。
\trFloer{}的思想以及进一步的相关发展将在续篇中详细讲述。
第6章尤其表明，源自\trCelestialMechanics{}的古老问题仍会引出强有力的新技术，它们对数学的其他分支同样有用。
应当指出，对一般\trSymplecticManifold{}的\trArnold{}猜想，在维数 $\geq 8$ 时仍未解决。

附录收录了一些技术性专题，以方便读者。
在A.1中我们证明：\trSymplectic\trDiffeomorphism{}可以局部地借一个单独的函数表示，即所谓的\textbf{\trGeneratingFunction}（generating function）。
这一经典事实在第5章中用到。
附录A.2以第4章中出现的\textbf{\trIntegrableSystem}（integrable system）的\trActionAngle{}坐标构造为例说明\trGeneratingFunction。
分析\trActionFunctional{}（第5章）所需的一个特殊\trSobolev\trEmbedding{}定理在A.3中证明。
我们推导了\trSphere{}上\trCauchyRiemann\trOperator{}的一些基本\textbf{\trEstimate}（estimate）（A.4）以及\trBoundary{}附近的\trEllipticEstimate{}（elliptic estimate）（A.5），并证明广义\trCarleman{}\trSimilarityPrinciple（A.6）；这些关于特殊\trPartialDifferentialEquation{}的结果全部在第6章中发挥重要作用。
前五章所需的分析工具书中都作了详细介绍，而\trTopology{}工具我们则不加解释地使用：我们用到\trBrouwer{}映射度（第2章）、\trLeray-\trSchauder{}度（第3章）、\trSmale{}度（$\bmod\ 2$）以及（上）\trHomology{}理论（第6章）。
有关这些\trTopology{}专题的参考文献见A.7与A.8，在那里我们讲解了\trBrouwer{}度以及\trAlexander-\trSpanier\trCohomology{}的\textbf{\trContinuity}（continuity）。
这一\trContinuity{}对我们在一般情形证明\trArnold{}猜想十分重要。

\subsection*{致谢}

对许多有益的意见、想法与建议，以及对撰写并完成本书的鼓励，我们非常感激。
我们特别要感谢M.~\trChaperon、Y.~\trEliashberg、S.~\trKuksin、J.~\trMoser、D.~\trSalamon、J.~C.~\trSikorav、C.~\trViterbo{}与K.~\trWysocki。
我们深受M.~\trBialy{}与L.~\trPolterovich{}之惠，他们在结果发表之前便惠示了关于\trGeodesic{}的成果。
对于细心阅读手稿、发现众多错误、加以订正并改进表述的C.~\trAbbas、K.~\trCieliebak、M.~\trFlucher、A.~\trGoing、B.~\trGratza、L.~\trKaas、M.~\trKriener、M.~\trSchwarz{}与K.~F.~\trSiburg，我们深表谢意。
伯尔尼\trBorel{}讨论班参加者们的关注令我们欣喜，我们也感激他们提出的问题。
感谢R.~\trMcLachlan{}为改进我们的英语所给的宝贵帮助，感谢克劳迪娅·\trFlepp{}以高效与耐心一遍遍录入并重新录入手稿。

\medskip
\noindent
苏黎世联邦理工学院，1994年3月28日
